% Lab 1 --- compact version (full docs: docs/web/docs/labs/lab1/)
\section{Lab 1: Reconnaissance and Enumeration}

Lab~1 introduces the reconnaissance phase of a penetration test following the PTES
framework~\cite{ptes}. Students interact with a fictional target organisation,
\textit{PEBCAK Corp}, through a layered topology that includes a perimeter firewall
(pfSense~CE~\cite{pfsense}), an internal router (VyOS~\cite{vyos}), and a target
Ubuntu~Server~24.04~\cite{ubuntu}. Access is exclusively via DNS hostname
(\texttt{lab1}) resolved through pfSense's Unbound resolver~\cite{unbound_dns}, requiring students
to use proper service enumeration before any exploitation attempt. A teacher reference --- credentials, network table, and startup checklist --- is compiled
in Appendix~B; the complete topology, configuration files, and step-by-step guide are at
\path{docs/web/docs/labs/lab1/}.


\begin{figure}[H]
\centering
\includegraphics[width=0.95\textwidth]{lab1-topology.png}
\caption{Lab~1 network topology}
\label{fig:lab1-topology}
\end{figure}

\newpage
\subsubsection*{Network Addressing}

\begin{table}[H]
\centering
\resizebox{\textwidth}{!}{\begin{tabular}{@{}lllll@{}}
\toprule
\textbf{Node} & \textbf{Role} & \textbf{Interface} & \textbf{IP} & \textbf{Network} \\
\midrule
Parrot-OS   & Attacker        & eth0 & 192.168.0.x (DHCP) & 192.168.0.0/24 (Homelab)  \\
pfSense     & Firewall/NAT    & WAN  & 192.168.0.29       & 192.168.0.0/24            \\
pfSense     & Firewall/NAT    & LAN  & 172.16.0.1         & 172.16.0.0/30             \\
VyOS        & Router/L3       & eth0 & 172.16.0.2         & 172.16.0.0/30             \\
VyOS        & Router/L3       & eth6 & 192.168.10.5       & 192.168.10.0/24 (Users)   \\
VyOS        & Router/L3       & eth7 & 192.168.20.1       & 192.168.20.0/24 (Servers) \\
Server      & Target (web+SSH)& eth0 & 192.168.20.50      & 192.168.20.0/24           \\
PC1         & Workstation     & eth0 & 192.168.10.50      & 192.168.10.0/24           \\
\bottomrule
\end{tabular}}
\caption{Lab~1 IP addressing scheme}
\label{tab:lab1-ip}
\end{table}

\subsubsection*{Techniques and Framework Alignment}

\begin{table}[H]
\centering
\small
\begin{tabularx}{\textwidth}{lXl}
\toprule
\textbf{Technique} & \textbf{Description} & \textbf{Reference} \\
\midrule
DNS resolution        & Hostname-to-IP via Unbound resolver (RFC~1035~\cite{rfc1035}) & PTES Recon  \\
Port scanning         & TCP SYN scan with Nmap (RFC~793~\cite{rfc793})                & PTES Recon  \\
Service enumeration   & Banner grabbing, version detection                            & PTES Recon  \\
Web content discovery & Directory brute-force with Gobuster                           & OWASP WSTG  \\
Credential retrieval  & Exposed credentials in web content                            & OWASP A05:2021\\
SSH access            & Password-based login (RFC~4252~\cite{rfc4252})                & PTES Exploit\\
\bottomrule
\end{tabularx}
\caption{Lab~1 techniques and framework mapping}
\label{tab:lab1-techniques}
\end{table}

\subsection{Problems Encountered and Solutions}

\subsubsection{VNC Keyboard Input on QEMU Nodes}

EVE-NG provides VNC console access to QEMU nodes, but certain key combinations ---
most notably \texttt{Ctrl+C} and special characters such as \texttt{|}, \texttt{@},
and \texttt{\{} --- are not reliably forwarded by the browser-based VNC client or by
the \texttt{vncdo} command-line tool. The root cause was that vncdo sent key symbols
rather than raw key codes, and the VNC server's key mapping did not cover all
combinations used in a penetration testing context.

The solution was to inject keystrokes directly via the QEMU monitor socket,
using Python's raw socket interface to send \texttt{sendkey} commands. This bypasses
the VNC protocol entirely and ensures that arbitrary key sequences, including those
involving modifier keys, are delivered correctly to the guest.

\subsubsection{pfSense DNS Resolver WAN Interface}

After initial configuration, DNS queries from the attacker machine were not being
resolved by pfSense's Unbound resolver. The problem was that Unbound, by default,
only listens on the LAN interface, so queries arriving on the WAN-side interface
were silently discarded. The fix was to explicitly add the WAN interface to
Unbound's listening configuration under \textit{Services $\to$ DNS Resolver $\to$
Network Interfaces}.

\subsubsection{Unbound Startup Configuration Path}

On subsequent restarts, the Unbound configuration applied via the
GUI was not being preserved, causing DNS resolution to fail after every reboot.
Investigation revealed that pfSense's startup scripts write the Unbound configuration
to \path{/var/unbound/unbound.conf} at boot time, overwriting any manual edits.
The correct approach is to configure all resolver settings exclusively through the
GUI or through pfSense's \texttt{/etc/inc/} PHP hooks, never by editing the
generated configuration file directly.

\subsubsection{VyOS Source NAT Requirement}

Traffic routed through VyOS from the attacker to the target was being dropped
because the target's default gateway was pfSense, not VyOS. Return packets were
therefore routed back via pfSense rather than through VyOS, breaking the connection.
The fix was to configure Source NAT (masquerade) on VyOS for traffic entering
the internal segment, ensuring that return traffic would follow the same path back
through the VyOS router.

\subsubsection*{Results}

\begin{table}[H]
\centering
\resizebox{\textwidth}{!}{\begin{tabular}{@{}lll@{}}
\toprule
\textbf{Flag} & \textbf{Location} & \textbf{Technique} \\
\midrule
\texttt{FLAG\{p3bc4k\_s3rv3r\_0wn3d\}} & \texttt{\textasciitilde/flag.txt} on Server via SSH & Credentials from hidden web page \\
\texttt{FLAG\{d3f4ult\_cr3ds\_k1ll\_s3cur1ty\}} & \texttt{/root/flag.txt} on pfSense via SSH pivot & Default credentials \\
\texttt{FLAG\{sk1n3t\_b3c4m3\_s3nt13nt\}} & \texttt{\textasciitilde/flag.txt} on PC1 via lateral move & Credentials found in pfSense flag \\
\bottomrule
\end{tabular}}
\caption{Lab~1 flags and acquisition techniques}
\label{tab:lab1-flags}
\end{table}
